% Supplementary B — Journal §II (Related Work)
% Full literature taxonomy with thematic clustering.

\documentclass[11pt,a4paper]{article}
% Shared preamble for all supplementary chapter documents.
% Used by each supplementary/conference/suppl_*.tex and
% supplementary/journal/suppl_*.tex via % Shared preamble for all supplementary chapter documents.
% Used by each supplementary/conference/suppl_*.tex and
% supplementary/journal/suppl_*.tex via % Shared preamble for all supplementary chapter documents.
% Used by each supplementary/conference/suppl_*.tex and
% supplementary/journal/suppl_*.tex via \input{../preamble.tex}.

\usepackage[utf8]{inputenc}
\usepackage[T1]{fontenc}
\usepackage[a4paper,margin=2.2cm]{geometry}
\usepackage{amsmath,amssymb,amsfonts,amsthm}
\usepackage{mathtools}
\usepackage{graphicx}
\usepackage{booktabs}
\usepackage{array}
\usepackage{multirow}
\usepackage{longtable}
\usepackage{algorithm}
\usepackage{algpseudocode}
\usepackage{xcolor}
\usepackage{url}
\usepackage{cite}
\usepackage[hidelinks]{hyperref}

\graphicspath{{../../figures/}}

\renewcommand{\thesection}{\Alph{section}.\arabic{section}}
\renewcommand{\thesubsection}{\thesection.\arabic{subsection}}

\theoremstyle{plain}
\newtheorem{theorem}{Theorem}[section]
\newtheorem{lemma}[theorem]{Lemma}
\newtheorem{proposition}[theorem]{Proposition}
\theoremstyle{definition}
\newtheorem{definition}[theorem]{Definition}
\theoremstyle{remark}
\newtheorem{remark}[theorem]{Remark}

\setcounter{secnumdepth}{3}
\setcounter{tocdepth}{2}

\newcommand{\R}{\mathbb{R}}
\newcommand{\E}{\mathbb{E}}
\newcommand{\KL}[2]{\mathrm{KL}\!\left(#1 \,\middle\|\, #2\right)}
\newcommand{\ari}{\mathrm{ARI}}
\newcommand{\nmi}{\mathrm{NMI}}
\newcommand{\softmax}{\mathrm{softmax}}
.

\usepackage[utf8]{inputenc}
\usepackage[T1]{fontenc}
\usepackage[a4paper,margin=2.2cm]{geometry}
\usepackage{amsmath,amssymb,amsfonts,amsthm}
\usepackage{mathtools}
\usepackage{graphicx}
\usepackage{booktabs}
\usepackage{array}
\usepackage{multirow}
\usepackage{longtable}
\usepackage{algorithm}
\usepackage{algpseudocode}
\usepackage{xcolor}
\usepackage{url}
\usepackage{cite}
\usepackage[hidelinks]{hyperref}

\graphicspath{{../../figures/}}

\renewcommand{\thesection}{\Alph{section}.\arabic{section}}
\renewcommand{\thesubsection}{\thesection.\arabic{subsection}}

\theoremstyle{plain}
\newtheorem{theorem}{Theorem}[section]
\newtheorem{lemma}[theorem]{Lemma}
\newtheorem{proposition}[theorem]{Proposition}
\theoremstyle{definition}
\newtheorem{definition}[theorem]{Definition}
\theoremstyle{remark}
\newtheorem{remark}[theorem]{Remark}

\setcounter{secnumdepth}{3}
\setcounter{tocdepth}{2}

\newcommand{\R}{\mathbb{R}}
\newcommand{\E}{\mathbb{E}}
\newcommand{\KL}[2]{\mathrm{KL}\!\left(#1 \,\middle\|\, #2\right)}
\newcommand{\ari}{\mathrm{ARI}}
\newcommand{\nmi}{\mathrm{NMI}}
\newcommand{\softmax}{\mathrm{softmax}}
.

\usepackage[utf8]{inputenc}
\usepackage[T1]{fontenc}
\usepackage[a4paper,margin=2.2cm]{geometry}
\usepackage{amsmath,amssymb,amsfonts,amsthm}
\usepackage{mathtools}
\usepackage{graphicx}
\usepackage{booktabs}
\usepackage{array}
\usepackage{multirow}
\usepackage{longtable}
\usepackage{algorithm}
\usepackage{algpseudocode}
\usepackage{xcolor}
\usepackage{url}
\usepackage{cite}
\usepackage[hidelinks]{hyperref}

\graphicspath{{../../figures/}}

\renewcommand{\thesection}{\Alph{section}.\arabic{section}}
\renewcommand{\thesubsection}{\thesection.\arabic{subsection}}

\theoremstyle{plain}
\newtheorem{theorem}{Theorem}[section]
\newtheorem{lemma}[theorem]{Lemma}
\newtheorem{proposition}[theorem]{Proposition}
\theoremstyle{definition}
\newtheorem{definition}[theorem]{Definition}
\theoremstyle{remark}
\newtheorem{remark}[theorem]{Remark}

\setcounter{secnumdepth}{3}
\setcounter{tocdepth}{2}

\newcommand{\R}{\mathbb{R}}
\newcommand{\E}{\mathbb{E}}
\newcommand{\KL}[2]{\mathrm{KL}\!\left(#1 \,\middle\|\, #2\right)}
\newcommand{\ari}{\mathrm{ARI}}
\newcommand{\nmi}{\mathrm{NMI}}
\newcommand{\softmax}{\mathrm{softmax}}


\title{Journal Supplementary B --- Full related-work taxonomy\\
{\large Companion to: \emph{Beyond Accuracy} (Section II)}}
\author{Felipe Santib\'a\~nez-Leal}
\date{2026}

\begin{document}
\maketitle

\section{Scope}
This supplement expands the related-work review of journal Section II
with thematic clustering and direct method-level comparison. We
organise the literature into four bands: (i)~foundational topic
modelling, (ii)~neural topic models, (iii)~HSI classification and
unmixing including topic-flavoured work, and (iv)~the evaluation
methodology literature this paper draws on.

\section{Foundational topic modelling}

\paragraph{LDA and its inference variants.} Blei, Ng, and
Jordan~\cite{blei2003lda} introduced the canonical Dirichlet
prior-over-topic-mixtures generative model with variational EM for
inference. Griffiths and Steyvers~\cite{griffiths2004finding}
developed the collapsed Gibbs alternative which avoids the
variational distribution but has higher per-iteration cost.
Hoffman, Blei, and Bach~\cite{hoffman2010online} introduced
online variational Bayes with natural-gradient ascent on the global
Dirichlet pseudo-counts, enabling LDA fits on corpora that do not
fit in memory; the scikit-learn implementation
\cite{pedregosa2011scikit} we use in this paper follows the
Hoffman 2010 recipe.

\paragraph{Non-parametric extensions.} Teh, Jordan, Beal, and
Blei~\cite{teh2006hdp} introduced the Hierarchical Dirichlet
Process as a non-parametric prior that lets the data choose the
topic count. We do not use HDP here because the band-frequency
tokenisation produces relatively short documents on which the HDP's
truncation-level estimate is unstable, and the experimental
budget would have multiplied without proportional information gain.

\paragraph{Topic interpretability tooling.} Sievert and Shirley
\cite{sievert2014ldavis} introduced LDAvis, the canonical
interactive interpretation tool for topic-word distributions, with
the relevance score $\lambda \phi_{k,b} + (1-\lambda) \log
\phi_{k,b}/p_b$. We use $\lambda = 0.5$ for the per-topic top-word
extraction. Taddy~\cite{taddy2012estimation} provided the formal
relationship between topic-word and word-topic probabilities used
in the journal's B-7 axis.

\section{Neural topic models}

\paragraph{ProdLDA.} Srivastava and Sutton~\cite{srivastava2017autoencoding}
introduced ProdLDA by reformulating LDA in a product-of-experts
form and combining it with the auto-encoding variational Bayes
framework of Kingma and Welling~\cite{kingma2014auto}. The Laplace
approximation to the Dirichlet prior is the technical core that
makes amortised inference tractable.

\paragraph{ETM.} Dieng, Ruiz, and Blei~\cite{dieng2020topic}
introduced the Embedded Topic Model, which decomposes the topic-word
distribution into an $V \times E$ word-embedding matrix and a
$K \times E$ topic-embedding matrix and softmax-normalises their
inner product. The low-rank decomposition lets rare words inherit
the topic associations of their semantic neighbours.

\paragraph{Neural variational topic models.} Miao, Yu, and Blunsom
\cite{miao2016neural} introduced the broader family of neural
variational document models (NVDM) of which ProdLDA can be seen as
a constrained variant; Higgins et al.~\cite{higgins2017betavae}
provided the $\beta$-VAE perspective that the deep-baseline ladder
in the companion code repository draws on.

\paragraph{What is missing.} The published neural-topic-model
literature for HSI is thin. ProdLDA and ETM have been applied to
text corpora extensively but to HSI rarely; the small body of HSI
papers that cite either method typically applies them on a single
scene with a single seed. The seed-stability swarms we report
(axis B-3, journal Figure 4) are, to our knowledge, the first
multi-seed ProdLDA and ETM evaluations on the six labelled
benchmarks.

\section{HSI classification, unmixing, and topic-flavoured work}

\paragraph{Linear unmixing.} Nascimento and Bioucas-Dias
\cite{nascimento2005vca} introduced the VCA endmember extraction
algorithm. The broader unmixing literature is reviewed by
Bioucas-Dias et al.~\cite{bioucasdias2012unmixing}. Linear unmixing
shares with LDA the non-negative basis structure but differs in two
operative ways: (a)~unmixing assumes a fixed endmember count
determined by physical reasoning, LDA admits $K$ as a free
hyperparameter; (b)~unmixing produces a per-pixel abundance vector
that is constrained to sum to one across endmembers, LDA produces a
per-document Dirichlet posterior whose mean satisfies the same sum
constraint asymptotically but admits a richer family of finite-data
shapes.

\paragraph{Deep spectral--spatial classification.} Plaza et al.
\cite{plaza2009recent} reviewed the first decade of HSI
classification, Camps-Valls et al.~\cite{camps2014advances} the
following half-decade with explicit attention to manifold and
kernel methods. Chen et al.~\cite{chen2014deep} introduced the
first systematic deep architectures for HSI classification and
Paoletti et al.~\cite{paoletti2019deep} surveyed the
subsequent five years. The OA-as-summary discipline in this
literature is well-developed and our journal paper does not
challenge it; the challenge is that topic-model adaptations
inherited the discipline without inheriting the inductive bias that
makes it well-suited to the supervised classification task.

\paragraph{Topic models on HSI.} Zou and Zare~\cite{zou2017pmlda}
introduced partial-membership LDA on a small set of HSI benchmarks.
Liu et al.~\cite{liu2019hsilda} explored a spectral-derivative
tokenisation. Both papers focus on improving the classification OA
of the topic representation; neither reports the cross-axis
diagnostics that the journal paper's twelve-axis framework
recommends.

\section{Evaluation methodology}

\paragraph{Topic-model coherence.} Röder, Both, and Hinneburg
\cite{roder2015exploring} surveyed coherence measures and showed
that the sliding-window $c_v$ measure tracks human judgements of
interpretability more reliably than perplexity. We use $c_v$,
$c_{\text{NPMI}}$, and U-Mass for axis B-2.

\paragraph{Clustering comparison.} Hubert and Arabie
\cite{hubert1985ari} formalised the adjusted Rand index. Vinh,
Epps, and Bailey~\cite{vinh2010information} extended the ARI
toolbox with the AMI and normalised information variants and
discussed when each is preferable.

\paragraph{Bipartite matching.} Kuhn~\cite{kuhn1955hungarian}
introduced the Hungarian algorithm for the maximum-weight bipartite
matching problem that the topic-id permutation step (axis B-8)
solves. The numerical implementation we use is in
SciPy~\cite{virtanen2020scipy,scipy_assignment}.

\paragraph{Reproducibility discipline.} Pineau et al.
\cite{pineau2021nature} argued for ML reproducibility checklists.
The journal paper's commitment to a single canonical fit + per-axis
deterministic builder + JSON artefact + smoke test is one
instantiation of that checklist.

\paragraph{Information-theoretic complements.} Koltchinskii
\cite{koltchinskii2000empirical} provides the empirical-process
foundation for the stability bounds used in seed-stability
analysis. Achanta et al.~\cite{achanta2012slic} introduced SLIC,
which we use as a cross-method-agreement baseline in axis B-6;
Felzenszwalb and Huttenlocher~\cite{felzenszwalb2004efficient}
provided the spatial-graph segmentation we use as the
non-spectral counterpart.

\section{Where this paper sits in the taxonomy}

The journal paper does not introduce a new topic model, a new
tokenisation, or a new inference procedure. It contributes a
\emph{multi-axis evaluation framework} that subsumes the
above-cited evaluation tools (Röder coherence, Hubert--Arabie ARI,
Kuhn matching, scikit-learn LDA, SciPy hypothesis tests) into a
single artefact pipeline that yields a per-axis numeric per
benchmark scene plus a hierarchical-Bayesian summary across scenes.
The contribution is methodological rather than algorithmic.

\bibliographystyle{IEEEtran}
\bibliography{../../bibliography/refs}

\end{document}
